\section{The actual original model from the auxiliary host construction}
\label{sec:cubic-auxiliary-actual-model}

The auxiliary form has a larger rational level, but it is used only to
produce support on a finite host.  The points subsequently placed on
that support are restrictions of the original fixed-type model.
This distinction permits the support construction to be nonempty
while the original Type II locus is empty.  We verify the complete
composition, including the two different support ideals and the
comparison with the original coefficient ring.

\paragraph{Original domain and retained source inputs.}
Fix
\[
 p=269,\quad K=\mathbb Q(\sqrt{-3}),\quad N=4,\quad
 \psi=\chi_{-4},\quad e\equiv267\pmod {268},\quad e>0,\quad
 j\in\{1,2\}.
 \tag{AM.1}
\]
Write \(c_p\) for the action on roots of unity, so
\(c_p(\operatorname{Fr}^{\rm geom}_w)=q_w^{-1}\), and put
\(\gamma=\psi c_p^e\).
Keep the original coefficient DVR \(\mathcal O_0\), fraction
field \(E_0\), residue field \(k_0\), and actual normalized
cubic ring \(S_j\) of AC-G65--AC-G66.  Its model is
\[
 (A_0,\mathcal A,x,q_0:A_0\longrightarrow S_j,\mathcal Q),
 \qquad I_0=\mathcal Q.\mathrm{ideal}.
 \tag{AM.2}
\]
Here \(q_0\) in (AM.2) is a ring map; the rational Eisenstein
auxiliary prime will always be written \(q_{\mathrm{aux}}=3229\).
Let \(\mathscr B_0\) be the complete degree-zero subring of
\(S_j\), and \(L_j\) the actual weight-minus-one cotangent
length.  None of these original objects is replaced by a
host deformation ring.

Retain the literal original \(\mathrm{GD}/\mathrm{Dat}\), marked
elements, coefficient, selected potentially crystalline component,
and local hypotheses of AC-G65--AC-G68.  Retain
\(\mathrm{CohInp}\), \(\mathrm{LemmaP}\), the generic algebraic
\(\mathrm{LemmaM}\) and \(\mathrm{LemmaM'}\), and the old
\(\mathrm{AutomorphyUnramifiedCoefficientsAlt}\) input used to
exclude original Type II points.  In particular impose
\[
                   e_{\mathcal O_0}
       =\operatorname{length}_{\mathcal O_0}(\mathcal O_0/p)=1
 \tag{AM.3}
\]
where required by that original modularity input.
The published attachment, ordinary control, transfer, local--global
compatibility and cyclic base-change inputs in AC-G119--AC-G120,
the native \(\mathrm{LevEx}\), \(\mathrm{ProModularRingExists}\),
\(\mathrm{LamEx}\), \(\mathrm{ModU}\), and the published
adjoint theorem of AC-G110 remain explicit inputs.
For the counting conclusion retain precisely the rank-three
original-field descent, witness and \(\mathrm{Hypotheses'}\)
inputs of AC-G111 and AC-G112-count.
The class-field, duality and Ferrero--Washington inputs retain
their scopes in AC-G74, AC-G90--AC-G100.

No original rational Type II endpoint is assumed to exist.
No mixed Type II Selmer vanishing or
\(\mathrm{RigiditySecondCase}\) hypothesis is used in this
composition.

\begin{theorem}[AC-G123-actual-model: actual support and the original integral bound]
\label{AC-G123-actual-model}
Under the preceding retained inputs, the auxiliary construction
supplies a finite CM host \(H/F\), \(F=H^+\), its actual full local
groups, a finite host coefficient DVR \(\mathcal O\) with residue
field \(k\), and the tuple rings
\[
 B\subset A\subset R,\qquad
 \theta=(\bar\chi^j\bar\gamma,1,\bar\chi^j\bar\gamma^{-1}).
 \tag{AM.4}
\]
For fixed \(\mathcal O_0\) and \(j\), these host data and
support ideals may be chosen once for the whole progression
\(e\equiv267\pmod {268}\).  Each original model still retains
its own stated original hypotheses.
For the actual native modular ideal
\(J_3\subset\mathcal B_3\), and the actual map
\(\Phi_3:\mathcal B_3\to B\), put \(I=\Phi_3(J_3)B\).
Then
\[
 I\subset\mathfrak q\cap B
       \quad\text{for every Type I prime }\mathfrak q\subset A.
 \tag{AM.5}
\]
This is the precise point-level support condition (G110.2).
There is also an ideal \(I_{\mathrm{off}}\subset B\) with
\[
 \sqrt{I_{\mathrm{off}}+\lambda_\theta}=\mathfrak m_B,\qquad
 I_{\mathrm{off}}\subset\mathfrak q\cap B
       \quad\text{for every off-triangular prime }\mathfrak q.
 \tag{AM.6}
\]
Here \(\lambda_\theta\) is the actual joint-residual-kernel
weight ideal.  It is the same ideal used in AC-G109, with
the host uniformizer included.

The original ring satisfies
\[
                   \mathscr B_0
          \text{ is finite over }\mathcal O_0.
 \tag{AM.7}
\]
Every original two-sided characteristic-zero local point is
Type I and rigid.  Consequently the original model (AM.2)
satisfies all five literal conditions \(\mathrm{M1}\)--\(\mathrm{M5}\).

For the exact original strict group
\[
 \begin{split}
 \alpha_j&=\chi^{-j}\psi^{-1}c_p^{-e}|_{G_K},\\
 \mathcal S_j&=\ker\Bigl(
  H^1(G_K,(E_0/\mathcal O_0)(\alpha_j))\\
 &\qquad\longrightarrow
  \prod_{v\mid p}H^1(G_{K_v},(E_0/\mathcal O_0)(\alpha_j))\\
 &\qquad\qquad\times\prod_{v\nmid p}H^1(I_v,(E_0/\mathcal O_0)(\alpha_j))
                         \Bigr),\\
 s_j&=\operatorname{length}_{\mathcal O_0}(\mathcal S_j)/
                                      e_{\mathcal O_0},
 \end{split}
 \tag{AM.8}
\]
the lengths are finite, and the conclusions of AC-G122-count
hold without an additional actual-support or \(\mathrm{M2}\)
premise.  Explicitly, take geometric Frobenius \(\sigma\) at
the degree-two place of \(K\) above \(11\), and put
\[
             b=121^e+1+121^{-e},\qquad Y=4(e+1).
 \tag{AM.9}
\]
There is a finite family \(X_j\) of actual Type I representations,
pairwise nonisomorphic on \(G_K\), such that
\[
 s_j\le \frac{L_j}{e_{\mathcal O_0}}
 \le 2\sum_{\rho\in X_j}
               v_{269}(\operatorname{tr}\rho(\sigma)-b),
 \qquad \#X_j\le C_1Y^{A_1}.
 \tag{AM.10}
\]
The constants \(C_1,A_1\) are exactly those of AC-G112-count.
For every nonempty set of labels \(J_0\subset\{1,2\}\), with
the retained hypotheses for those original models, there is
one nonzero integer \(\mathcal N^I_{11,J_0}\) such that
\begin{align}
 \sum_{j\in J_0}L_j/e_{\mathcal O_0}
     &\le 2v_{269}(\mathcal N^I_{11,J_0}),\notag\\
 \log|\mathcal N^I_{11,J_0}|
     &\le 5C_1\#J_0\,Y^{A_1+1}\log11,\notag\\
 \left(\sum_{j\in J_0}L_j/e_{\mathcal O_0}\right)\log269
     &\le10C_1\#J_0\,Y^{A_1+1}\log11.
 \tag{AM.11}
\end{align}
The same bounds hold with \(s_j\) in place of \(L_j/e_{\mathcal O_0}\).
The supports and hosts for the two labels may be chosen separately;
the common integer uses their one common original-field classifier.
There is no change of the original coefficient ring or valuation.
\end{theorem}

\begin{proof}
\proofstep{1. A nonempty auxiliary starting point and the exact initial host.}
Set \(e_*=267\).  AC-G119 constructs the actual ordinary
newform of weight \(2e_*+1=535\) and its representation \(V\), using
\(q_{\mathrm{aux}}=3229\).  The order-\(p\) character is
removed on its totally real cyclic field \(T_\theta\).
AC-G120 supplies the actual seed and fields
\[
 \begin{gathered}
 F_1=K\mathbb Q(\sqrt{-23}),\quad E=\mathbb Q(\sqrt5),\\
 H_{\mathrm{ab}}=F_1E,\qquad
 H_0=F_1T_\theta D,\qquad F_0=H_0^+,\\
 H_{\mathrm{ab}}\subset H_{\mathrm{ab}}T_\theta\subset H_0,
                     \quad\text{both indices }p .
 \end{gathered}
 \tag{AM.12}
\]
Here \(D/E\) is the new cyclic ray extension of degree \(p\).
The two auxiliary fields are not confused: \(T_\theta\) kills
the Eisenstein type before the unitary construction; \(D\)
kills the new level-raising type afterward.
The rank-three seed over \(H_0\) is absolutely irreducible
and ordinary of the symmetric weight for \(e_*\).
It has \(N=0\), away inertia \((\psi,1,\psi^{-1})\), and
the exact ordered \(p\)-flag with finite tuple
\((\chi^j,1,\chi^j)\) and powers \(c_p^{e_*},1,c_p^{-e_*}\).
The rank-two representation \(V|_{G_{H_0}}\) is absolutely
irreducible and has the corresponding rational two-character
flag and away type.

The extension \(H_0/F_0\) is finite-unramified and relatively
split at \(p,2,3\).  Its \(p\)-completions are the unramified
degree-\(538\) fields.  Its intersection with the original
prime-to-\(p\) residual splitting field \(\Omega\) is \(K\).
These are the actual statements proved in AC-G120; no
original \(\mathrm{Dat.IsPt}\) is supplied or used.
For the class estimate use only the independent finite
arithmetic certificate on \(H_{\mathrm{ab}}\) in AC-G103.
It gives adjacent cyclotomic rank at most \(3\), against
real degree \(4\).  AC-G100 at the two steps in (AM.12)
therefore gives
\[
                 r_{\alpha,H_0}\le \tfrac34[F_0:\mathbb Q],
 \qquad \alpha=\bar\chi^{-j}\bar\gamma^{-1}.
 \tag{AM.13}
\]
Each step is unramified at \(p\) and disjoint from the
totally ramified cyclotomic tower.  This is the precise
normalized rank transport; it does not apply
Ferrero--Washington to the cubic splitting field.
Since \(\bar c_p\) has order \(268\), the actual residual
character \(\bar\psi\bar c_p^{\,e}\) equals the one for
\(e_*\) for every weight in (AM.1).
All the class sectors in this construction are consequently
the same sectors throughout that progression.

\proofstep{2. A common continuation supplies filling and the strict dimension gap.}
Apply AC-G92 to \(F_0\), with the full prime-to-\(p\)
residual splitting extension \(M/F_0\), containing
\(H_0\), \(\mu_p\), and all the displayed residual characters.
Choose at least one selected place.  Its cyclic extension
\(T'/F_0\) has degree \(p^s\) and splits every \(p\)-place.
It need not be Galois over \(\mathbb Q\).
The host \(\widetilde H=H_0T'\) remains CM and relatively
finite-unramified; its \(p\)-completions are unchanged.
All relevant fields are actual fields in that construction.

The exact class quotient of AC-G93 supplies a nonzero strict
adjacent class and a nonzero central anti class.  It also
supplies the rational sectors
\(\bar\gamma^{-1}\eta_K\) and \(\bar\gamma^{-2}\eta_K\)
needed in AC-G121.  Their oddness makes them nontrivial.
At \(p\), the cubic order three cannot cancel a rational
inertia character of order dividing \(268\).  The endpoint
square has order \(134\).  Thus all the local invariant
vanishings used to identify the strict cone groups hold.
They persist after \(p\)-power restriction.

For clarity, the bounded dual used to fill the adjacent
class is exactly
\[
       \lambda=\eta_K\bar\psi^2\bar c_p^{\,2e+1}.
 \tag{AM.14}
\]
It is even over \(\mathbb Q\), its splitting field is
totally real and abelian, and its \(p\)-inertia is nontrivial:
\(2e+1\equiv267\pmod {268}\).
After adjoining that abelian splitting field and \(\mu_p\),
the real host is a finite \(p\)-extension of an abelian
field.  The extra selected extension is again cyclic of
\(p\)-power degree.  This is exactly the finite-\(p\)
ascent and full-local duality situation proved in
AC-G93-fill, including its non-\(\mathbb Q\)-Galois case.
The dual \(H^1\) is bounded; its eventually isomorphic
restriction maps have zero corestriction modulo \(p\).
Perfect duality kills the mixed strict \(H^2\) obstruction
after one further layer.  Nontriviality of the residual
characters makes restriction injective on the chosen
adjacent and central classes.  Therefore the actual
normalized long-entry filling and Low path persist on
all sufficiently high layers.

The selected \(p\)-split extension is disjoint from the
cyclotomic tower, so AC-G100 also transports (AM.13) to
\(\widetilde H\).  Reindex after eventual total ramification
and write \(H_n/F_n\) for its cyclotomic layers and
\(d_n=[F_n:\mathbb Q]=d_0p^n\).  The exact AC-G98 formula is
\[
 a_n=r_\alpha p^n+\sum_i\min(t_i,p^n),\qquad
 h_n=O(1),\quad z_n=O(1),\qquad r_\alpha\le\tfrac34d_0 .
 \tag{AM.15}
\]
The adjacent permutation correction is zero by its nontrivial
local cubic character.  The bounded central and trivial terms
involve only rational coefficient fields and finite-\(p\)
ascent from abelian fields.  Hence, on every sufficiently
high layer,
\[
                    h_n+a_n+z_n+1<d_n,\qquad d_{v,n}>5 .
 \tag{AM.16}
\]
Reindexing multiplies both \(r_\alpha\) and \(d_0\) by the
same power of \(p\), so it does not invalidate the strict
ratio in (AM.15).
Choose one layer past all these thresholds and the filling
threshold.  Call it \(H/F\).  The rational-frame thresholds
in AC-G121 can be met at the same time.  No effective bound
on this layer, and no cubic \(\mu=0\), is being assumed.

\proofstep{3. The actual rank-three native ideal and strong seed.}
The rank-three seed restricts irreducibly through the
selected and cyclotomic steps: a nontrivial cyclic
quotient-character self-twist would have cube one by
determinants, whereas each step has degree a power of \(p\).
The retained cyclic base-change theorem preserves its
cuspidal automorphy.  Its finite away type and ordered
ordinary \(p\)-flag restrict to the full local groups.
The seed weight is still \(e_*=267\); there is no
need to replace this seed by an unverified low-weight form.

Choose the finite host coefficient DVR \(\mathcal O\)
to contain the coefficients, local-character values and
all \(p\)-adic embeddings needed for both actual seeds
and native \(\mathrm{Large}\).  AC-G121 constructs the
split residual rank-two lattice after a finite value
extension if necessary.  This choice is made before
either host universal space is formed and does not
alter \(\mathcal O_0\).
The rank-three native determinant only needs its actual
residual determinant, not a claim of split matrix reduction
for an arbitrary rank-three stable lattice.

The normal closure over \(H_0\) of the selected extension
and any cyclotomic layer is a \(p\)-extension.
Its intersection with \(\Omega\) is still \(K\).
The scalar residual element with
\(\chi=1,\bar\gamma=1,\bar c_p=-1\) therefore extends
fixing this normal closure.  A fresh completely split
rational auxiliary place gives the actual \(\mathrm{V0ok}\).
The determinant-only argument of AC-G89 gives \(\mathrm{N4}\)
for every native point.  The host has no \(\mu_p\);
its normal-closure degree has only the prime factors \(2,p\).
The exact \(\mathrm{HidaSetting}\), \(\mathrm{Lev}\),
full local maps and away joint kernel are thus those
constructed in AC-G121 and AC-G97.  The source
\(\mathrm{p\_unr}\) concerns this fresh auxiliary place.

Use the actual tuple space of AC-G94--AC-G95.
The arithmetic Euler and finiteness fields are AC-G96,
the generators and class exponent are AC-G99, and the
residual local separation is unchanged by these \(p\)-steps.
AC-G97 gives its native Pair and
\(\Phi_3:\mathcal B_3\to B\); AC-G102 gives surjectivity
and finite weight fiber for the image of a native modular
ideal.  These are the changed-tuple constructions, not
the original single-character rank-three \(\mathrm{TwLink}\).

The actual seed \(r_H\), multiplied by \(c_p^{-1}\),
defines the native determinant/local-character point.
Indeed compactness supplies an integral stable lattice after
the chosen finite value extension.  Its determinant law is
integral and has the required residual semisimplification.
The ordered local characters are integral units and reduce
to their labelled residual characters.  Set
\(\nu=r_Hc_p^{-1}\), take \(D=\det(\nu)\), and take each
native local character to be \(c_p^{-1}\) times the
corresponding character of the full flag.
The actual conjugate-self-dual polarization gives
\(\mathrm{Pair.P}\), and the full ordered flags give
\(\mathrm{Pair.Ord}\).  If \(g\) lies in the away joint
kernel, then \(r_H(g)=I\) and \(c_p(g)=1\).
Consequently \(\nu(g)=I\), and the determinant trace
of \([h]([g]-1)\) is zero for every \(h\), exactly
as required by \(\mathrm{Pair.I}\).
The local characters kill the closed commutator; the host
has no local \(p\)-power roots of unity, so the stipulated
\(\mathrm{mu}/\mathrm{th}\) condition is exactly the one
in AC-G121.  These checks give an actual local map from
the native universal ring to \(\mathcal O\).
Its away joint kernel is killed exactly, since its
finite prime-to-\(p\) inertia is
\((\psi,1,\psi^{-1})\); there is no residual unipotent
or monodromy term.  Its full local flags are the
AC-G120 flags after restriction.
Writing the native weights at a parallel exponent \(w\) as
\[
                         h(w)=(1-w,1,1+w),
 \tag{AM.17}
\]
the seed has weights \(h(e_*)\), which increase strictly.
No pair of its local characters
differs by \(c_p\): on a finite-index subgroup killing
the finite characters this would give a nonzero power
among \(c_p^{-1},c_p^{\pm e_*-1},c_p^{\pm2e_*-1}\) equal
to one.  All exponents are nonzero here and local
cyclotomic inertia is infinite.  It is unramified at
the fresh auxiliary place.
Thus every clause of native \(\mathrm{PT}\) holds.
The retained \(\mathrm{ModU}\) supplies a proper ideal
\(J_3\) killed by this actual seed point.

Its native quotient is nonzero finite torsionfree over
\(\mathcal O[[X_1,\ldots,X_{3d}]]\), and it has
\(\mathrm{Unr0}\) and the exact native prime-core properties.
The absolutely irreducible seed realizes its determinant.
These are all premises of (G107.4).
AC-G107 now recovers a whole native component, with the
actual identity \(\beta\Phi_3=q\), and gives
\[
 \xi\in\operatorname{Spec}A\text{ of Type I},\quad
 Z=\xi\cap B,\quad I\subset Z,\quad \dim B/Z=3d+1.
 \tag{AM.18}
\]
No torsionfreeness of \(B/I\) has been inferred from that
of the native quotient.  This is the precise strong
seed \((S)\), not just a characteristic-zero point.
The full ordinary functor and its native residual datum use
\(\bar\gamma\), the fixed scalar twist \(c_p^{-1}\), and
the full local maps.  The local \(\mathrm{mu}\) is the
closed commutator and \(\mathrm{th}=1\).  They impose no
fixed characteristic-zero Hodge weight.
Thus the rings, \(J_3\), the rational support of AC-G121,
and the weight ideal are the same for every \(e\) in (AM.1).
The single seed at \(e_*\) supplies their strong component;
it is not replaced by a conjectural original-weight seed.
This is independence of the reduced tuple and native
ordinary datum, not equality of the full integral local
characters.  The finite cubic and quadratic inertia
characters are fixed, and the Teichmuller component of
\(c_p^e\) is fixed on this progression, whereas its
principal-unit powers do vary.  Each original model has
its own fixed integral flag with powers \(c_p^{\pm e}\).
The native full ordinary functor permits these varying
powers; in Step 5 each characteristic-\(p\) curve reduces
its own integral flag to the one common residual tuple.

\proofstep{4. Supply both support assertions with their correct ideals.}
AC-G121 gives one uniform even ideal \(M_E\subset B\):
\[
 \sqrt{M_E+\lambda_\theta}=\mathfrak m_B,\qquad
 M_E\subset P\cap B
       \quad\text{for every even off-triangular prime }P.
 \tag{AM.19}
\]
It applies to every such prime, with the middle-character
deformation retained.  This is condition \((E)\).
The actual native ideal image \(I=\Phi_3(J_3)B\) separately
has finite weight fiber by AC-G102.
Its normalized prime-core properties are those of AC-G104:
they use the actual native \(\mathrm{P1core}/\mathrm{P3core}\)
and the \(\mathrm{N4}\) just supplied.
The local presentations, normalized collision bounds and
component-chain ideal are AC-G101 and AC-G105.
Their full-local Euler hypotheses are AC-G96, and their
numerical inequalities are (AM.16).
The filled adjacent and central classes of Step 2 give
the actual connected punctured lower core of AC-G95.
Thus every premise of AC-G108 is now supplied on this host.

Applying that theorem with \(J=I\), \(J_E=M_E\), and
the strong seed (AM.18) proves (AM.5), and proves (AM.6)
with
\[
                         I_{\mathrm{off}}=I\cap M_E .
 \tag{AM.20}
\]
The two ideals are deliberately distinguished.
It is \(I=\Phi_3(J_3)B\), not an arbitrary finite-weight
ideal, which lies below every Type I prime.  Therefore
(AM.5) is exactly the required native condition (G110.2).
It is \(I_{\mathrm{off}}\) which handles both types for
the integral characteristic-\(p\) curve argument.

\proofstep{5. Compare every original curve and every original Type I point.}
Take any complete-DVR characteristic-\(p\) curve
\(f:S_j\to V\) in the literal original integral finiteness
test.  Fix compatible embeddings of \(k_0,k\) into an
algebraic closure of \(\mathbb F_p\).
If \(V=k'[[t]]\), choose finite \(k''\) containing \(k'\)
and \(k\), inject \(V\) into \(V'=k''[[t]]\), and use
the complete host test algebra
\[
                          C=k+tV'.
 \tag{AM.21}
\]
It carries \(\mathcal O\to k\) and has residue field exactly
\(k\).  All entries of the restricted representation and
its inverse lie in \(C\), since their reductions are the
prescribed tuple values.  No scalar extension or quotient
identification of the original ring is made.

The original actual away-kernel killing condition kills
every generator of the host joint kernel after restriction.
At \(p\), AC-G65-flag fixes the integral inertia characters
before the curve is chosen.  They become their constant
residual characters on this characteristic-\(p\) curve.
The labelled Hensel projectors descend the full ordered
flags to \(C\); the strict \(J\)-orthogonal normalization
at the host witness preserves them.  The actual AC-G94
functor gives \(R\to C\), and on the joint residual
inertia kernel each local character is one.
The specialization therefore kills \(\lambda_\theta\),
exactly as in AC-G109-curve.

If this restricted curve were off both triangular loci,
(AM.6) would make its \(B\)-image the closed residue point.
A surviving positive entry \(x\) of weight \(a>0\) and
negative entry \(y\) of weight \(-b<0\) would then give
a nonzero image of the zero-weight monomial
\(x^b y^a\in\mathfrak m_B\), a contradiction.
The host curve is triangular.  Its diagonal characters,
after removing their residual characters, are unramified
everywhere and have finite class-group exponent.
The principal units of a characteristic-\(p\) domain
have no torsion, so these remaining characters are one.
Consequently the host characteristic polynomials are
the constant seed polynomials.

For any \(g\in G_K\), some positive power of \(g\) lies
in \(G_H\).  Its eigenvalues therefore satisfy fixed
finite-field equations.  The characteristic coefficients
are algebraic over the finite residue field and have the
prescribed reductions, hence are precisely those constants.
This reasoning permits \(p\) to divide the index.
The generic shape and matrix trace-pairing argument of
AC-G109 excludes the irreducible and endpoint alternatives
for this curve.  Thus its original coordinates have all
weight-zero values zero and one entire nonzero-weight
sign zero.  The integral \(\mathrm{finCrit}\) theorem,
applied over \(\mathcal O_0\), gives (AM.7).
This verifies the original \(\mathrm{M2}\), including
vertical components, rather than only generic-fiber
finiteness.

For a characteristic-zero Type I local point of \(S_j\),
choose a finite value field containing its matrices and
the host coefficient field.  In its ring of integers use
the inverse image of \(k\) as the complete host test ring,
as in (G110.5).  Restriction, exact away-kernel killing,
the fixed local flags and the same strict orthogonal
normalization give an actual field point of \(A\).
AC-G110-irr proves that its restriction to \(G_H\) and
to \(G_{H(\mu_{p^\infty})}\) is irreducible.  Its kernel
is therefore a Type I prime.
Equation (AM.5) makes the composite
\(\mathcal B_3\xrightarrow{\Phi_3}B\to\overline{\mathbb Q}_p\)
kill the actual \(J_3\).
The determinant and full local characters are exactly
those of the original restricted representation multiplied
by \(c_p^{-1}\); its own weights are \(h(e)\) in (AM.17).
The \(\mathrm{end2}\) direction of native classicality
now gives its actual automorphy.
This is the point-level implication of (G110.2), not
an invocation of \(\mathrm{PT}\) in the reverse direction.
The polarized adjoint theorem and restriction/corestriction
then give the original rational adjoint vanishing of
AC-G110-rigidity.  The coefficient field has characteristic
zero, so even a \(p\)-divisible host degree is invertible
in this cohomological descent.

\proofstep{6. Original rigidity and the unchanged integral counting normalization.}
The original shape theorem makes every two-sided local point
Type I or Type II.  AC-G116, or its quadratic extension
AC-G117, excludes every actual Type II point on (AM.1)
under the retained original modularity and coefficient
hypotheses.  It excludes a point of \(S_j\), not the
auxiliary seed on \(H\).  Hence all original two-sided
points are Type I and are rigid by Step 5.
The exact model orbit/dual-point extension and normalized
coboundary argument of AC-G122-M4 gives \(\mathrm{M4}\)
over \(\mathcal O_0\).  In detail, normalize any two
point derivations by the same grading-torus element,
extend them to actual dual points, and use rigidity.
Their normalized cobounding matrices lie in the one
line spanned by \(\operatorname{diag}(1,0,-1)\).
Entry determination then makes the two derivations
linearly dependent; undoing the nonzero weight scalars
is the literal \(\mathrm{M4}\) quantifier.
No mixed Type II group is declared zero.

AC-G66 already supplies \(\mathrm{M1},\mathrm{M3}\),
and Step 5 supplies \(\mathrm{M2}\).
For \(\mathrm{M5}\), the homogeneous original trace
difference at \(11\) has value zero at the origin and
value \(\operatorname{tr}\rho(\sigma)-b\) at every
two-sided point.  The seed value is \(b\), since both
the rational quadratic and cubic characters are trivial
at the degree-two Frobenius.  It is strictly larger than
three.  The actual Type I trace has complex absolute
value at most three, so the difference is nonzero.
The torus orbit extension to every model point fixes
this degree-zero element.  This is AC-G113-M5 with
the Type II part unnecessary.

The host has a tower from \(K\) consisting of cyclic
steps of degree \(2\) or \(p\): the two quadratic
factors are already in \(H_{\mathrm{ab}}\), followed
by the two \(p\)-steps, the selected extension and
the cyclotomic layers.
Thus the exact original-field descent of AC-G111
applies to every original Type I point now made
automorphic.  Its original inertia conditions, not the
auxiliary seed level, give the principal compact level
\(269N\) of AC-G112.  The primes \(q_{\mathrm{aux}}\)
and the ray prime do not enter that original conductor.
The exact original-field \(\mathrm{WitA}\), unlinking,
and \(\mathrm{Af.TrData}\) bound-three statements are
those already proved there with their retained sources.

All hypotheses of AC-G122-count have now been supplied.
It applies \(\mathrm{LemmaM}\) with
\(v_{\mathcal O_0}=e_{\mathcal O_0}v_{269}\), obtains
the orbit-reduced set \(X_j\), and uses the actual
\(\mathrm{LemmaM'}\) strict injection over the same
\(\mathcal O_0\).  This proves (AM.8)--(AM.10).
Only the rank-three classifier is needed because every
actual point is Type I.  Its coefficient-independent
norm construction uses the stable label set \(J_0\)
with \(z(j)=1\), including two distinct labels with
the same value when \(J_0=\{1,2\}\).
The resulting exact bounds are (AM.11), with
constant \(20C_1\) in the last inequality for both labels.
No rescaling of length by the host coefficient ramification
index appears anywhere.
\end{proof}

\paragraph{Scope of the completed composition.}
This theorem supplies the actual support premises of AC-G122
on (AM.1); it does not add them as new unexplained assumptions.
All stated native, published and generic algebraic inputs remain.
The conclusion concerns the original fixed-type integral models
at \(269\), and includes their cotangent lengths.
AC-G114 already proves the strict groups in (AM.8) zero by
a separate finite certificate; that fact is not used to infer
model rigidity or support here.
There is no all-character or all-coefficient-prime estimate,
no treatment of omitted primes \(2,3\), and no conclusion
about the numerical family-average budget or H04.4.
The generic singular regulator problem for other Type II
families is not removed.

The finite-cover interpretation is the previously audited
SGA~1, Expose~V, sections~4 and~7 dictionary.
The auxiliary cover kills its own finite character; subsequent
restriction and normalization place the original representation
on a different, actual host chart.  The native ideal propagates
through that chart and is then tested on original curves and
points.  No identification of the original deformation ring
with the host ring, no new motive, and no fresh archival theorem
is asserted.
